\documentclass[11pt,compress,t,notes=noshow, xcolor=table]{beamer}

\input{../../style/preamble}
\input{../../latex-math/basic-math}
\input{../../latex-math/basic-ml}
\input{../../latex-math/optim-basics}

\title{Optimization in Machine Learning}

\begin{document}

\titlemeta{
Second order methods
}{
Newton-Raphson
}{
figure/NR_quad_approx_good.png
}{
\item Newton-Raphson
\item Damped Newton-Raphson
\item Convergence of Newton-Raphson
\vspace{0.5cm}
\item[] \hspace{-0.7cm}Slides based on\\
\hspace{-0.7cm}\furtherreading{AGGARWAL2020} (Ch. 5.4-5.7)
}

\begin{framei}{Motivation: Beyond gradient information}
\item \textbf{Recall:} GD only uses the gradient $\Rightarrow$ \textit{ignores curvature}
\item Hessian $\H = \nabla^2 f$ (curvature) encodes how fast the gradient itself changes along each direction
\begin{itemizeM}
\item High curvature: gradient is unreliable further away $\Rightarrow$ take a \textit{small} step
\item Low curvature: gradient stays valid further away $\Rightarrow$ take a \textit{large} step
\end{itemizeM}
\item GD only picks one global $\alpha$ for all directions $\Rightarrow$ for ill-conditioned $\H$, this causes zig-zagging / slow convergence:
\imageC[0.55]{../04-multivariate-first-order/figure/curvature_2.png}
\end{framei}


\begin{framei}[sep=L]{Motivation: Beyond gradient information}
\item \textbf{Idea:} Use full curvature information instead of a single scalar $\alpha$
\begin{itemizeM}
\item rescale the gradient \textit{individually per direction} via $\H^{-1}$
\item biases updates towards low-curvature directions, where a large step is unlikely to overshoot or worsen $f$
\end{itemizeM}
\imageC[0.35][https://link.springer.com/book/10.1007/978-3-030-40344-7]{figure_man/aggarwal_fig5-7_inverse_hessian.jpg}
\begin{center}{\footnotesize Effect of multiplying the steepest-descent direction with the inverse
Hessian}\end{center}
\end{framei}


\comment[NOTE]{BB}{vielleicht noch eine "demotivierende" slide warum 2nd order schlöecht ist in hohen dim
oder kommt das in quasi newton?
wir müssen das wirklich mehr in deatils analysieren das ist wichtig. am besten mit bezug auf grosse daten und speed vergleoich. vielleicht sollte das ein extra chunk sein}

\begin{framei}{Newton-Raphson}
\item \textbf{Assumption:} $f \in \CC{2}$
\item \textbf{Aim:} Find stationary point $\xv^\ast$, i.e., $\nabla f(\xv^\ast) = \zero$
\item \textbf{Idea:} Find root of first order Taylor approx of $\nabla f(\xv)$:
$$
\nabla f(\xv) \approx \nabla f(\xv^{[t]}) +
\nabla^2 f(\xv^{[t]})(\xv - \xv^{[t]}) = \zero
$$
$$
\nabla^2 f(\xv^{[t]})(\xv - \xv^{[t]}) = - \nabla f(\xv^{[t]})
$$
$$
\xv^{[t+1]} = \xv^{[t]} - (\nabla^2 f(\xv^{[t]}))^{-1}\nabla f(\xv^{[t]})
$$
\item \textbf{Update scheme:}
$$
\xv^{[t+1]} = \xv^{[t]} + \mathbf{d}^{[t]}
$$
with $\mathbf{d}^{[t]} = - (\nabla^2 f(\xv^{[t]}))^{-1}\nabla f(\xv^{[t]})$
\end{framei}


\begin{framei}{Newton-Raphson: Practice}
\item \textbf{Note:} In practice, we get $\mathbf{d}^{[t]}$ by solving the linear system
$$
\nabla^2 f(\xv^{[t]})\mathbf{d}^{[t]} = - \nabla f(\xv^{[t]})
$$
with direct (matrix decompositions) or iterative methods
\end{framei}


\begin{framei}{Newton-Raphson: Intuition}
\item Newton update approximates the function locally with a quadratic bowl and moves to the bottom
of the bowl in a single step
\item The ``learning rate'' is incorporated implicitly
\item In general, the bottom of the quadratic bowl is not the bottom of the true function $\Rightarrow$
multiple Newton updates are usually necessary
\imageC[0.55]{figure/NR_quad_approx_good.png}
\item[] {\small \textbf{Note}: univariate functions are rather ``boring'' as there are only two possible directions $\Rightarrow$ we'll move on to more interesting examples soon}
\item \textbf{Caveat}: if the quadratic approximation is poor (e.g., far from the current point), the Newton step may worsen $f$ (see later examples)
\end{framei}


\begin{framei}{Analytical example with quadratic form}
\item Consider $f(x_1, x_2) = x_1^2 + \frac{x_2^2}{2}$
\item Update direction: $\mathbf{d}^{[t]} = -(\nabla^2 f(x_1^{[t]}, x_2^{[t]}))^{-1} \nabla f(x_1^{[t]}, x_2^{[t]})$
\item Gradient and Hessian:
$$
\nabla f(x_1, x_2) = \begin{pmatrix}2x_1 \\ x_2\end{pmatrix}, \quad
\nabla^2 f(x_1, x_2) = \begin{pmatrix}2 & 0 \\ 0 & 1\end{pmatrix}
$$
\item First step:
$$
\begin{pmatrix}
x_1^{[1]} \\ x_2^{[1]}
\end{pmatrix} = \begin{pmatrix}
x_1^{[0]} \\ x_2^{[0]}
\end{pmatrix} + \mathbf{d}^{[0]} = \begin{pmatrix}
x_1^{[0]} \\ x_2^{[0]}
\end{pmatrix} - \begin{pmatrix}
1/2 &  0 \\ 0 & 1
\end{pmatrix} \begin{pmatrix}
2x_1^{[0]} \\ x_2^{[0]}
\end{pmatrix}
$$
$$
= \begin{pmatrix}
x_1^{[0]} \\ x_2^{[0]}
\end{pmatrix} + \begin{pmatrix}
-x_1^{[0]} \\
-x_2^{[0]}
\end{pmatrix} = \zero
$$
\end{framei}


\begin{framev}{Analytical example with quadratic form}
\splitV[0.55]{
\begin{itemizeM}
\item Newton-Raphson only needs one iteration for quadratic forms
\item See how the inverse Hessian rescales the gradient:
\begin{itemizeM}
\item $x_1$: higher curvature $\Rightarrow$ scaled by $1/2$ by $\H^{-1}$
\item $x_2$: lower curvature $\Rightarrow$ scaled by $1$ by $\H^{-1}$
\end{itemizeM}
\item \textbf{Note}: For diagonal $H$, each dimension is scaled individually and we can easily interpret the scaling factors. For non-diagonal $H$, the update direction becomes a linear combination of all dimensions.
\end{itemizeM}
}{
\imageC[1]{figure/NR_quadratic_form.png}
}
\end{framev}


\begin{framev}{Non-quadratic examples: Convergence}
\textbf{Example 1}: Newton-Raphson on
$f(\xv) = \exp(x_1 + 3 x_2 - 0.1) + \exp(x_1 - 3 x_2 - 0.1) + \exp(-x_1 - 0.1)$

\vspace{0.5cm}
\only<1>{
After $1$ Newton step:
\splitV[0.55]{
\imageC[1]{figure/NR_boyd_1_contour.png}
}{
\imageC[1]{figure/NR_boyd_1_error.png}
}
}
\only<2>{
After $2$ Newton steps:
\splitV[0.55]{
\imageC[1]{figure/NR_boyd_2_contour.png}
}{
\imageC[1]{figure/NR_boyd_2_error.png}
}
}
\only<3>{
After $5$ Newton steps $\Rightarrow$ Newton-Raphson converges in just 5 steps:
\splitV[0.55]{
\imageC[1]{figure/NR_boyd_5_contour.png}
}{
\imageC[1]{figure/NR_boyd_5_error.png}
}
}
\end{framev}


\begin{framev}{Non-quadratic examples: Overshooting}
\textbf{Example 2}: Newton-Raphson on
$f(\xv) = (1 + x_1^2)^{1/2} + \frac{x_2^2}{2}$\\
\textit{Note:} along $x_1$, curvature decays away from the optimum:
$\frac{\partial^2}{\partial x_1^2} (1 + x_1^2)^{1/2} = (1 + x_1^2)^{-3/2} \to 0$

\vspace{0.3cm}
\only<1>{
After $1$ Newton step:
\splitV[0.55]{
\imageC[1]{figure/NR_overshoot_1_contour.png}
}{
\imageC[1]{figure/NR_overshoot_1_error.png}
}
$\Rightarrow$ quadratic in $x_2$, direction solved exactly in one step
}
\only<2>{
After $2$ Newton steps:
\splitV[0.55]{
\imageC[1]{figure/NR_overshoot_2_contour.png}
}{
\imageC[1]{figure/NR_overshoot_2_error.png}
}
$\Rightarrow$ along $x_1$, the quadratic approx. is too flat, its minimum lies well beyond $\xv^\ast$ $\Rightarrow$ NR overshoots
}
\only<3>{
After $3$ Newton steps:
\splitV[0.55]{
\imageC[1]{figure/NR_overshoot_3_contour.png}
}{
\imageC[1]{figure/NR_overshoot_3_error.png}
}
$\Rightarrow$ NR diverges \\
\vspace{0.3cm}
\textit{Note}: $\H \succ 0$ everywhere $\Rightarrow$ direction is fine, only step size is wrong
}
\end{framev}


\begin{framei}{Newton-Raphson: Step size}
\item \textbf{Relaxed/Damped Newton-Raphson:} Use step size $\alpha^{[t]} > 0$ with
$$
\xv^{[t+1]} = \xv^{[t]} + \alpha^{[t]} \mathbf{d}^{[t]}
$$
to satisfy Wolfe conditions (or just Armijo rule)
\item \textbf{Example 2 (cont'd)}: Same function and starting point, now with Armijo backtracking
($\alpha^{[t]} \leftarrow 1$, halved until $f$ decreases sufficiently):
\splitV[0.55]{
\imageC[1]{figure/NR_overshoot_damped_contour.png}
}{
\imageC[1]{figure/NR_overshoot_damped_error.png}
}
\vfill
Only iteration $2$ is damped ($\alpha = 0.5$) $\Rightarrow$ enough for NR to converge
\end{framei}


\begin{framei}[sep=L]{Discussion}
\item \textbf{Advantage:}
\item[] For $f$ sufficiently smooth:
\begin{framed}
\centering
Newton-Raphson converges \textit{locally} quadratically \\
(i.e., for starting points close enough to stationary point) \\
$\Rightarrow$ \emph{see ``Deep Dive: Convergence of Newton-Raphson'' slides}
\end{framed}
\vfill
\item \textbf{Disadvantage:}
\item[] For \enquote{bad} starting points:
\begin{framed}
\centering
Newton-Raphson may diverge \\
$\Rightarrow$ \emph{see ``Limitations of Newton-Raphson'' slides}
\end{framed}
\end{framei}

\endlecture
\end{document}
